\documentclass[10pt]{extarticle}
\usepackage[a4paper,margin=1.5cm]{geometry}
\usepackage[T1]{fontenc}
\usepackage{amsmath,amssymb}
\usepackage{graphicx}
\usepackage[export]{adjustbox}
\usepackage{multicol}
\usepackage{enumitem}
\usepackage{booktabs,array,tabularx}
\setlength{\columnsep}{0.7cm}
\setlength{\parindent}{0pt}
\setlength{\parskip}{3pt}
\pagestyle{empty}
\makeatletter
% figure registry (figdefs.tex is generated by the extract step)
\newcommand{\qfigdef}[4]{\@namedef{qf@file@#1}{#2}\@namedef{qf@w@#1}{#3}\@namedef{qf@wide@#1}{#4}}
\newcommand{\qf@check}[1]{\@ifundefined{qf@file@#1}{\PackageError{score2}{Unknown figure id #1}{}}{}}
\newcommand{\qf@img}[4][]{\includegraphics[width=#2\dimexpr\@nameuse{qf@w@#4}pt\relax,max width=#3,max height=0.45\textheight,#1]{\@nameuse{qf@file@#4}}}
% Full-width questions: finalize wraps every question that contains a wide figure in \qfullbegin...\qfullend,
% so the WHOLE question (stem, figures, options) is set across the page in reading order, then the two columns
% resume. The question's head - its number, stem and first figure (\qfullsplit marks the end of the head) -
% is kept on one page, so a number is never left at a page foot with its figure overleaf; the rest may break
% across pages normally, so no big gaps are left. Without a \qfullsplit, the whole question is kept together
% if it is under 45% of a page. Never write these yourself.
\newif\ifqf@full
\newif\ifqf@boxed
\newsavebox{\qf@box}
\newcommand{\qf@open}{\global\qf@fulltrue\par\global\qf@boxedtrue
  \setbox\qf@box\vbox\bgroup\hsize\textwidth\linewidth\textwidth}
\newcommand{\qf@emit}[1]{\par\egroup\global\qf@boxedfalse
  \ifdim\dimexpr\ht\qf@box+\dp\qf@box\relax<#1\textheight\noindent\usebox{\qf@box}\par\else\unvbox\qf@box\fi}
\newcommand{\qfullsplit}{\ifqf@boxed\qf@emit{0.85}\fi}
\newcommand{\qf@close}{\ifqf@boxed\qf@emit{0.45}\fi\global\qf@fullfalse}
\newcommand{\qfullbegin}{\end{multicols}\qf@open}
\newcommand{\qfullend}{\qf@close\begin{multicols}{2}}
\newcommand{\qfullstart}{\qf@open}
\newcommand{\qfulldone}{\qf@close}
% \qfig[scale]{id}: block figure, centred, directly where it is written (size set by the tool).
% Fallback if a wide figure is not inside a full-width question: break out of the columns in place - no floats.
\newcommand{\qfig}[2][1]{\qf@check{#2}%
  \ifqf@full\qf@blk{#1}{#2}\else
  \ifnum\@nameuse{qf@wide@#2}=1 \expandafter\@firstoftwo\else\expandafter\@secondoftwo\fi
  {\end{multicols}\begin{center}\qf@img{#1}{\textwidth}{#2}\end{center}\begin{multicols}{2}}%
  {\qf@blk{#1}{#2}}\fi}
\newcommand{\qf@blk}[2]{\par\smallskip{\centering\qf@img{#1}{\linewidth}{#2}\par}\smallskip}
% \qfigi[scale]{id}: inline figure (inside option lists / tables / minipages); never breaks columns.
\newcommand{\qfigi}[2][1]{\qf@check{#2}\qf@img[valign=c]{#1}{\linewidth}{#2}}
\makeatother
% structure macros - lint reads these to check completeness and figure placement
\newcommand{\papersec}[1]{\par\bigskip{\centering\Large\bfseries PAPER-#1\par}\medskip}
\newcommand{\partsec}[2]{\par\medskip\noindent{\large\bfseries PART-#1: #2}\par\smallskip}
\newcommand{\sect}[1]{\par\smallskip\noindent\textbf{#1}\par}
\newcommand{\qn}[1]{\par\smallskip\noindent\textbf{#1.}~}
\newcommand{\paraq}[2]{\par\smallskip\noindent\textbf{Paragraph for Questions #1 and #2}\par}
\newenvironment{qoptions}{\begin{itemize}[leftmargin=1.8em,itemsep=1pt,topsep=2pt,label={}]}{\end{itemize}}
\newcommand{\opt}[2]{\item[(#1)] #2}
% figure registry (generated by score2_tools.py finalize)
\qfigdef{p1_phy_s1_q01_1}{TOPS - Test 03 (20 Nov 2022)-img/p1_phy_s1_q01_1.png}{91.8}{0}
\qfigdef{p1_phy_s1_q02_1}{TOPS - Test 03 (20 Nov 2022)-img/p1_phy_s1_q02_1.png}{99.9}{0}
\qfigdef{p1_phy_s1_q03_1}{TOPS - Test 03 (20 Nov 2022)-img/p1_phy_s1_q03_1.png}{143.7}{0}
\qfigdef{p1_phy_s1_q04_1}{TOPS - Test 03 (20 Nov 2022)-img/p1_phy_s1_q04_1.png}{69.4}{0}
\qfigdef{p1_phy_s2_q05_1}{TOPS - Test 03 (20 Nov 2022)-img/p1_phy_s2_q05_1.png}{100.6}{0}
\qfigdef{p1_phy_s2_q07_1}{TOPS - Test 03 (20 Nov 2022)-img/p1_phy_s2_q07_1.png}{122.7}{0}
\qfigdef{p1_phy_s2_q09_1}{TOPS - Test 03 (20 Nov 2022)-img/p1_phy_s2_q09_1.png}{178.0}{0}
\qfigdef{p1_phy_s2_q10_1}{TOPS - Test 03 (20 Nov 2022)-img/p1_phy_s2_q10_1.png}{100.4}{0}
\qfigdef{p1_phy_s2_q10_2}{TOPS - Test 03 (20 Nov 2022)-img/p1_phy_s2_q10_2.png}{101.5}{0}
\qfigdef{p1_phy_s3_para1_2_1}{TOPS - Test 03 (20 Nov 2022)-img/p1_phy_s3_para1_2_1.png}{110.6}{0}
\qfigdef{p1_phy_s3_para3_4_1}{TOPS - Test 03 (20 Nov 2022)-img/p1_phy_s3_para3_4_1.png}{192.7}{0}
\qfigdef{p1_phy_s3_para5_6_1}{TOPS - Test 03 (20 Nov 2022)-img/p1_phy_s3_para5_6_1.png}{145.2}{0}
\qfigdef{p1_phy_s4_q07_1}{TOPS - Test 03 (20 Nov 2022)-img/p1_phy_s4_q07_1.png}{172.5}{0}
\qfigdef{p1_phy_s4_q08_1}{TOPS - Test 03 (20 Nov 2022)-img/p1_phy_s4_q08_1.png}{109.7}{0}
\qfigdef{p1_phy_s4_q09_1}{TOPS - Test 03 (20 Nov 2022)-img/p1_phy_s4_q09_1.png}{82.5}{0}
\qfigdef{p1_chem_s1_q03_A}{TOPS - Test 03 (20 Nov 2022)-img/p1_chem_s1_q03_A.png}{31.8}{0}
\qfigdef{p1_chem_s1_q03_B}{TOPS - Test 03 (20 Nov 2022)-img/p1_chem_s1_q03_B.png}{24.0}{0}
\qfigdef{p1_chem_s1_q03_C}{TOPS - Test 03 (20 Nov 2022)-img/p1_chem_s1_q03_C.png}{66.1}{0}
\qfigdef{p1_chem_s1_q03_D}{TOPS - Test 03 (20 Nov 2022)-img/p1_chem_s1_q03_D.png}{24.0}{0}
\qfigdef{p1_chem_s2_q07_1}{TOPS - Test 03 (20 Nov 2022)-img/p1_chem_s2_q07_1.png}{21.5}{0}
\qfigdef{p1_chem_s2_q07_A}{TOPS - Test 03 (20 Nov 2022)-img/p1_chem_s2_q07_A.png}{28.9}{0}
\qfigdef{p1_chem_s2_q07_B}{TOPS - Test 03 (20 Nov 2022)-img/p1_chem_s2_q07_B.png}{74.5}{0}
\qfigdef{p1_chem_s2_q07_C}{TOPS - Test 03 (20 Nov 2022)-img/p1_chem_s2_q07_C.png}{80.4}{0}
\qfigdef{p1_chem_s2_q07_D}{TOPS - Test 03 (20 Nov 2022)-img/p1_chem_s2_q07_D.png}{78.5}{0}
\qfigdef{p1_chem_s2_q08_A}{TOPS - Test 03 (20 Nov 2022)-img/p1_chem_s2_q08_A.png}{125.8}{0}
\qfigdef{p1_chem_s2_q08_2}{TOPS - Test 03 (20 Nov 2022)-img/p1_chem_s2_q08_2.png}{21.5}{0}
\qfigdef{p1_chem_s2_q08_B}{TOPS - Test 03 (20 Nov 2022)-img/p1_chem_s2_q08_B.png}{18.7}{0}
\qfigdef{p1_chem_s2_q08_4}{TOPS - Test 03 (20 Nov 2022)-img/p1_chem_s2_q08_4.png}{26.7}{0}
\qfigdef{p1_chem_s2_q08_C}{TOPS - Test 03 (20 Nov 2022)-img/p1_chem_s2_q08_C.png}{26.5}{0}
\qfigdef{p1_chem_s2_q08_D}{TOPS - Test 03 (20 Nov 2022)-img/p1_chem_s2_q08_D.png}{135.7}{0}
\qfigdef{p1_chem_s2_q09_1}{TOPS - Test 03 (20 Nov 2022)-img/p1_chem_s2_q09_1.png}{62.0}{0}
\qfigdef{p1_chem_s3_para3_4_1}{TOPS - Test 03 (20 Nov 2022)-img/p1_chem_s3_para3_4_1.png}{150.6}{0}
\qfigdef{p1_chem_s4_q08_1}{TOPS - Test 03 (20 Nov 2022)-img/p1_chem_s4_q08_1.png}{60.4}{0}
\qfigdef{p1_math_s1_q04_1}{TOPS - Test 03 (20 Nov 2022)-img/p1_math_s1_q04_1.png}{63.0}{0}
\qfigdef{p1_math_s3_para5_6_1}{TOPS - Test 03 (20 Nov 2022)-img/p1_math_s3_para5_6_1.png}{121.0}{0}
\qfigdef{p2_phy_s1_q01_1}{TOPS - Test 03 (20 Nov 2022)-img/p2_phy_s1_q01_1.png}{134.5}{0}
\qfigdef{p2_phy_s1_q02_1}{TOPS - Test 03 (20 Nov 2022)-img/p2_phy_s1_q02_1.png}{57.6}{0}
\qfigdef{p2_phy_s1_q05_1}{TOPS - Test 03 (20 Nov 2022)-img/p2_phy_s1_q05_1.png}{63.2}{0}
\qfigdef{p2_phy_s1_q06_1}{TOPS - Test 03 (20 Nov 2022)-img/p2_phy_s1_q06_1.png}{171.1}{0}
\qfigdef{p2_phy_s2_q01_1}{TOPS - Test 03 (20 Nov 2022)-img/p2_phy_s2_q01_1.png}{192.0}{0}
\qfigdef{p2_phy_s2_q02_1}{TOPS - Test 03 (20 Nov 2022)-img/p2_phy_s2_q02_1.png}{39.7}{0}
\qfigdef{p2_phy_s2_q03_1}{TOPS - Test 03 (20 Nov 2022)-img/p2_phy_s2_q03_1.png}{139.8}{0}
\qfigdef{p2_phy_s2_q04_1}{TOPS - Test 03 (20 Nov 2022)-img/p2_phy_s2_q04_1.png}{149.6}{0}
\qfigdef{p2_phy_s2_q05_1}{TOPS - Test 03 (20 Nov 2022)-img/p2_phy_s2_q05_1.png}{141.3}{0}
\qfigdef{p2_phy_s2_q06_1}{TOPS - Test 03 (20 Nov 2022)-img/p2_phy_s2_q06_1.png}{93.0}{0}
\qfigdef{p2_phy_s3_q01_1}{TOPS - Test 03 (20 Nov 2022)-img/p2_phy_s3_q01_1.png}{120.5}{0}
\qfigdef{p2_phy_s3_q02_1}{TOPS - Test 03 (20 Nov 2022)-img/p2_phy_s3_q02_1.png}{154.6}{0}
\qfigdef{p2_phy_s3_q03_1}{TOPS - Test 03 (20 Nov 2022)-img/p2_phy_s3_q03_1.png}{36.2}{0}
\qfigdef{p2_phy_s3_q06_1}{TOPS - Test 03 (20 Nov 2022)-img/p2_phy_s3_q06_1.png}{124.9}{0}
\qfigdef{p2_chem_s1_q01_1}{TOPS - Test 03 (20 Nov 2022)-img/p2_chem_s1_q01_1.png}{108.0}{0}
\qfigdef{p2_chem_s1_q03_A}{TOPS - Test 03 (20 Nov 2022)-img/p2_chem_s1_q03_A.png}{26.7}{0}
\qfigdef{p2_chem_s1_q03_B}{TOPS - Test 03 (20 Nov 2022)-img/p2_chem_s1_q03_B.png}{24.0}{0}
\qfigdef{p2_chem_s1_q03_C}{TOPS - Test 03 (20 Nov 2022)-img/p2_chem_s1_q03_C.png}{26.7}{0}
\qfigdef{p2_chem_s1_q03_D}{TOPS - Test 03 (20 Nov 2022)-img/p2_chem_s1_q03_D.png}{24.0}{0}
\qfigdef{p2_chem_s2_q03_1}{TOPS - Test 03 (20 Nov 2022)-img/p2_chem_s2_q03_1.png}{202.7}{0}
\qfigdef{p2_chem_s2_q03_2}{TOPS - Test 03 (20 Nov 2022)-img/p2_chem_s2_q03_2.png}{212.8}{0}
\qfigdef{p2_chem_s2_q03_3}{TOPS - Test 03 (20 Nov 2022)-img/p2_chem_s2_q03_3.png}{155.6}{0}
\qfigdef{p2_chem_s2_q04_1}{TOPS - Test 03 (20 Nov 2022)-img/p2_chem_s2_q04_1.png}{225.5}{0}
\qfigdef{p2_chem_s2_q04_2}{TOPS - Test 03 (20 Nov 2022)-img/p2_chem_s2_q04_2.png}{231.7}{0}
\qfigdef{p2_chem_s2_q04_3}{TOPS - Test 03 (20 Nov 2022)-img/p2_chem_s2_q04_3.png}{262.8}{1}
\qfigdef{p2_chem_s3_q04_1}{TOPS - Test 03 (20 Nov 2022)-img/p2_chem_s3_q04_1.png}{81.6}{0}
\raggedcolumns
\makeatletter\@beginparpenalty=10000\makeatother
\newcolumntype{L}[1]{>{\raggedright\arraybackslash}p{#1}}
\begin{document}
\begin{center}{\Large\bfseries TOPS --- Test 03 (20 Nov 2022)}\end{center}
\begin{multicols}{2}
\raggedcolumns
\papersec{1}
\partsec{1}{PHYSICS}
\sect{SECTION-I(i)}
\qn{1} A uniform ruler of mass $m$ and length $\ell$ is balanced horizontally on the two index fingers as shown in the figure (top view). Initially the fingers are at the ends of the ruler. Now, one hand at a time, the index fingers are moved slowly in horizontal direction towards the other to meet at the middle of the ruler, during which the ruler slowly slides one way or the other. If the coefficients of static and kinetic frictions between the ruler and fingers are $\mu$ and $\mu/3$ respectively, then what is total work done by the external agent till the hands meet?
\qfig{p1_phy_s1_q01_1}
\begin{qoptions}
\opt{A}{$\dfrac{\mu mg\ell}4$}
\opt{B}{$\dfrac{\mu mg\ell}6\ln\left(\dfrac{3\sqrt3}2\right)$}
\opt{C}{$\dfrac{\mu mg\ell}2\ln\left(\dfrac52\right)$}
\opt{D}{None of these}
\end{qoptions}
\qn{2} A heavy disc with radius $R$ is rolling down hanging on two inextensible strings wound around the disc very tightly. The free ends of the strings are attached to a fixed horizontal support. The strings are always tensed during the motion. At some instant, the angular velocity of the disc is $\omega$, and the angle between the strings is $\alpha$. Find the velocity of centre of mass of the disc at this moment.
\qfig{p1_phy_s1_q02_1}
\begin{qoptions}
\opt{A}{$\dfrac{\omega R}{\sin\left(\frac\alpha2\right)}$}
\opt{B}{$\dfrac{\omega R}{\cos\left(\frac\alpha2\right)}$}
\opt{C}{$\dfrac{2\omega R}{\sin\left(\frac\alpha2\right)}$}
\opt{D}{$\dfrac{2\omega R}{\cos\left(\frac\alpha2\right)}$}
\end{qoptions}
\qn{3} A ball of mass $m=1$ kg is hung vertically by a thread of length $\ell=1.50$ m. Upper end of the thread is attached to the ceiling of a trolley of mass $M=4$ kg. Initially trolley is stationary and it is free to move along horizontal rails without friction. A shell of mass $m=1$ kg, moving horizontally with velocity $v_0=6$ ms$^{-1}$, collides with the ball and gets stuck with it. As a result, thread starts to deflect towards right. Calculate its maximum deflection with the vertical. ($g=10$ ms$^{-2}$)
\qfig{p1_phy_s1_q03_1}
\begin{qoptions}
\opt{A}{$30^\circ$}
\opt{B}{$37^\circ$}
\opt{C}{$53^\circ$}
\opt{D}{$60^\circ$}
\end{qoptions}
\qn{4} A small block $A$ of mass 1 kg is placed over a horizontal plate $B$ of mass 2 kg which is placed on a smooth horizontal surface. Block $A$ is being moved continuously 1 m/s as shown in the figure and the plate is moved continuously with speed 1 m/s, towards the direction shown in the figure. If the coefficient of friction between blocks is 0.1, the net power loss due to friction is : ($g=10$ m/s$^2$)
\qfig{p1_phy_s1_q04_1}
\begin{qoptions}
\opt{A}{0.5 J/s}
\opt{B}{1 J/s}
\opt{C}{$\sqrt2$ J/s}
\opt{D}{2 J/s}
\end{qoptions}
\sect{SECTION-I(ii)}
\qn{5} Two identical rods $OA$ and $AB$ each of mass $M$ and length $L$ are pinned to each other at joint $A$ and one of the rods is hinged to the smooth horizontal floor at $O$. Initially the system is released from rest when rod $OA$ makes an angle $\theta$ with the horizontal as shown in the figure. Which of the following options are correct when the pinned point $A$ is about to hit the floor ?
\qfig{p1_phy_s2_q05_1}
\begin{qoptions}
\opt{A}{Angular speed of the right rod is $\sqrt{\dfrac{3g\sin\theta}L}$.}
\opt{B}{Angular acceleration of the left rod is $\dfrac{3g}{2L}$.}
\opt{C}{Normal reaction between the right rod and floor is $\dfrac{Mg}8$.}
\opt{D}{None of the above.}
\end{qoptions}
\qn{6} A thin uniform rod is placed vertically on a horizontal rough floor. If it is disturbed by a negligible amount horizontally, which of the following statements are correct ?
\begin{qoptions}
\opt{A}{There exists a finite value of coefficient of friction for which the rod will not slip till it becomes horizontal.}
\opt{B}{For any finite value of coefficient of friction the rod will definitely slip before it rotates by angle of $\cos^{-1}\left(\dfrac13\right)$.}
\opt{C}{If it rotates by an angle of $60^\circ$ without slipping, the direction of friction on the rod has changed.}
\opt{D}{The rod might slip in the direction of fall or opposite depending on the value of coefficient of friction.}
\end{qoptions}
\qn{7} A uniform disc of mass $\eta m$ and radius $R$ is connected to a block of mass $m$ with a help of an inextensible ropes and massless spring of force constant $K$, as shown in the figure. If the pulleys are smooth and ground is rough enough to avoid slipping, answer the following when the system is released from rest with spring in natural state.
\qfig{p1_phy_s2_q07_1}
\begin{qoptions}
\opt{A}{If $\eta=1/8$, maximum angular acceleration of the disc is $\dfrac{2g}R$.}
\opt{B}{If $\eta=1/6$, maximum upwards acceleration of the block is $\dfrac{3g}5$.}
\opt{C}{If $\eta=1/2$, the coefficient of friction must be $\geqslant\dfrac4{13}$.}
\opt{D}{None of the above.}
\end{qoptions}
\qn{8} A rigid body is rotating about a fixed axis so that its angular velocity depends on total angle of rotation $\theta$ as $\omega=\omega_0-\alpha\theta$ where $\omega_0$ and $\alpha$ are positive constants and at time $t=0$, $\theta=0$. The correct option(s) of the following is/are :-
\begin{qoptions}
\opt{A}{$\theta$ at any instant is $\dfrac{\omega_0}\alpha\left(1-e^{-\alpha t}\right)$.}
\opt{B}{$\theta$ at any instant is $\dfrac{\omega_0}\alpha e^{-\alpha t}$.}
\opt{C}{Angular velocity at any instant is $\omega_0e^{-\alpha t}$}
\opt{D}{Number of rotations made by the body before coming to rest is $\dfrac{\omega_0}{2\pi\alpha}$.}
\end{qoptions}
\qn{9} A 13 ft ladder is leaning against the wall when its base starts to slide away. When at some instant of time the base is 12ft from the wall, then it was seen to be moving at a rate of 5 ft/sec. Select correct option(s) at the same instant of time
\qfig{p1_phy_s2_q09_1}
\begin{qoptions}
\opt{A}{The top of the ladder is sliding down the wall at a rate of magnitude 12 ft/s.}
\opt{B}{The area of the triangle formed by the ladder, wall and ground is changing at a rate of magnitude 60 ft$^2$/s.}
\opt{C}{The angle $\theta$ between the ladder and the ground is changing at a rate of magnitude 1 rad/s.}
\opt{D}{The angle $\theta$ between the ladder and the ground is changing at a rate of magnitude 4 rad/s.}
\end{qoptions}
\qn{10} Two balls are thrown from an inclined plane at angle of projection $\alpha$ with the plane one up the incline plane and other down the incline as shown in the figure. If $R_1$ \& $R_2$ be their respective ranges, then [here $T_1$ \& $T_2$ are times of flight in the two cases respectively]
\qfig{p1_phy_s2_q10_1}
\qfig{p1_phy_s2_q10_2}
\begin{qoptions}
\opt{A}{$h_1=h_2$}
\opt{B}{$R_2-R_1=T_1^2$}
\opt{C}{$R_2-R_1=g\sin\theta\,T_2^2$}
\opt{D}{$R_2-R_1=g\sin\theta\,T_1^2$}
\end{qoptions}
\sect{SECTION-II(i)}
\paraq{1}{2} A planar, uniform, thin symmetric plate of mass `$M$' and of shape as shown in the figure has moment of inertia $I$ about the axis $AA'$. If moment of inertia of this plate about the axis passing through the point `$O$' and perpendicular to plane of the plate is $I_1$ and also its moment of inertia about the axes $BB'$, $CC'$ and $DD'$ are $I_2$, $I_3$ and $I_4$ respectively. Given each side of triangle and square shown is of length $a$. Answer the following two questions.
\qfig{p1_phy_s3_para1_2_1}
\qn{1} The sum of moment of inertial terms $I_1$ and $I_2$ is $\left(3I+kMa^2\right)$. Find `$k$'.
\qn{2} The sum of moment of inertial terms $I_3$ and $I_4$ is $\left(2I+kMa^2\right)$. Find `$k$'
\paraq{3}{4} A bicycle has pedal rods of length 16cm connected to a spoked disc of radius 10cm. The bicycle wheels are 80cm in diameter and the chain runs over a gear of radius 4cm. The cyclist applies a force $F$ perpendicular to the pedal rod, but the cycle is rest. Assume that tension in the lower part of the chain is negligible. Neglect any other frictional force within the parts of the bicycle.
\qfig{p1_phy_s3_para3_4_1}
\qn{3} If the force applied by the person is $F=100$N but the bicycle is at rest. The sum of tension in the upper part of the chain and the frictional force on the rear wheel (in Hectonewton) is equal to \underline{\hspace{1cm}}.
\qn{4} When the person applies a force $F_1$ perpendicular to the pedal rod the bicycle is on the verge of moving and frictional force on the front wheel is 20N. The value of $F_1$ in Hectonewton is \underline{\hspace{1cm}}
\paraq{5}{6} Two blocks $A$ and $B$ of masses $m_1=2$kg and $m_2=8$ kg respectively are allowed to move without friction. Block $A$ is on block $B$ and block $B$ slides on smooth fixed inclined plane as shown in the figure. The angle of inclination of inclined plane is $45^\circ$. (Take $g=10$m/s$^2$)
\qfig{p1_phy_s3_para5_6_1}
\qn{5} The acceleration of the block $B$ with respect to the ground in m/s$^2$ is \underline{\hspace{1cm}}
\qn{6} The acceleration of the block $A$ with respect to the ground in m/s$^2$ is \underline{\hspace{1cm}}
\sect{SECTION-II(ii)}
\qn{7} A stick $AB$ of mass $M$ is tied to one of a light string $OA$. A horizontal force $F=Mg$ is applied at end $B$ of the stick and it remains in equilibrium in the position shown. If $\alpha=\tan^{-1}(n)$, the value of $n$ is
\qfig{p1_phy_s4_q07_1}
\qn{8} In the figure shown, $\mu_1=0.5$, $\mu_2=0.2$, $\mu_3=0.2$, $\mu_4=0.05$, between the pair of surfaces, from top to bottom. Find the acceleration of third block of mass 30kg.
\qfig{p1_phy_s4_q08_1}
\qn{9} Figure shows a hollow vertical cylinder of mass `$M$', radius `$R$' and height `$R$', kept on a smooth horizontal floor. It has smooth helical groove on the inner surface. Number of turns in the helical groove is large, so that any jerk can be neglected, when a particle of mass `$M$' enters the groove through a small hole near the bottom in the cylinder moving along the tangent on the floor. Take $R=10$cm and $g=10$m/s$^2$.
\qfig{p1_phy_s4_q09_1}
The minimum speed $v_0$ of the particle on the horizontal plane, so that it reaches the top of the cylinder is $\sqrt{KRg}$. Find $K$.
\columnbreak
\partsec{2}{CHEMISTRY}
\sect{SECTION-I(i)}
\qn{1} 20 ml of mixture of CH$_4$ and CO$_2$ when reacted with just sufficient oxygen, total volume after the complete reaction becomes 40 ml at 1 atm and 500K (for all volume measurement). The mass percentage of CH$_4$ in the initial mixture is :
\begin{qoptions}
\opt{A}{50}
\opt{B}{40}
\opt{C}{26.66}
\opt{D}{60}
\end{qoptions}
\qn{2} 0.518 g sample of lime-stone is dissolved in HCl and entire calcium is precipitated as CaC$_2$O$_4$. After filtering and washing the precipitate 40 ml of 0.25 N acidified KMnO$_4$ were required to titrate the precipitate in a redox titration. Calculate mass of CaO obtained from strong heating of the initial sample :
\begin{qoptions}
\opt{A}{0.13 g}
\opt{B}{0.56 g}
\opt{C}{0.26 g}
\opt{D}{0.28 g}
\end{qoptions}
\qn{3} Compound in which inductive, hyper conjugation and resonance all three effect are present?
\begin{qoptions}
\opt{A}{\qfigi{p1_chem_s1_q03_A}}
\opt{B}{\qfigi{p1_chem_s1_q03_B}}
\opt{C}{\qfigi{p1_chem_s1_q03_C}}
\opt{D}{\qfigi{p1_chem_s1_q03_D}}
\end{qoptions}
\qn{4} If there are three possible values of spin quantum number $\left\{s=+\dfrac12,0,-\dfrac12\right\}$ then maximum number of elements can be placed in 3$^{\text{rd}}$ period.
\begin{qoptions}
\opt{A}{8}
\opt{B}{18}
\opt{C}{12}
\opt{D}{27}
\end{qoptions}
\sect{SECTION-I(ii)}
\qn{5} One mole of KMnO$_4$ in acidic condition is used for complete oxidation of FeSO$_4$, FeC$_2$O$_4$ and H$_2$C$_2$O$_4$ respectively and separately. Pick up the correct statement :
\begin{qoptions}
\opt{A}{5 mole of FeSO$_4$ can be oxidised}
\opt{B}{$\dfrac35$ mole of FeC$_2$O$_4$ can be oxidised}
\opt{C}{$\dfrac53$ mole of FeC$_2$O$_4$ can be oxidised}
\opt{D}{2.5 mole of H$_2$C$_2$O$_4$ can be oxidised}
\end{qoptions}
\qn{6} A 4 M aqueous solution of H$_2$SO$_4$ has density 0.86 g/ml. Which is/are correct for this solution :
\begin{qoptions}
\opt{A}{Molality of solution is about 8.55 m.}
\opt{B}{62.5 ml of this solution required 50 ml of 40\% w/v NaOH solution for complete neutralisation.}
\opt{C}{If 30 moles of SO$_3$ is added to 1 liter of this solution then \% labelling of this solution is 102.2\%.}
\opt{D}{If 100 ml 2 M H$_3$PO$_4$ is added to 100 ml of this solution then concentration of H$^+$ in the final solution is 7 M. (Assuming complete dissociation of acids).}
\end{qoptions}
\qn{7} Identify correct acidic strength order?
\begin{qoptions}
\opt{A}{\qfigi{p1_chem_s2_q07_1}$<$\qfigi{p1_chem_s2_q07_A}}
\opt{B}{\qfigi{p1_chem_s2_q07_B}}
\opt{C}{\qfigi{p1_chem_s2_q07_C}}
\opt{D}{\qfigi{p1_chem_s2_q07_D}}
\end{qoptions}
\qn{8} Identify correct basic strength order :
\begin{qoptions}
\opt{A}{\qfigi{p1_chem_s2_q08_A}}
\opt{B}{\qfigi{p1_chem_s2_q08_2}$<$\qfigi{p1_chem_s2_q08_B}}
\opt{C}{\qfigi{p1_chem_s2_q08_4}$>$\qfigi{p1_chem_s2_q08_C}}
\opt{D}{\qfigi{p1_chem_s2_q08_D}}
\end{qoptions}
\qn{9} Select the incorrect option(s) regarding the orbital given in following diagram.
\qfig{p1_chem_s2_q09_1}
\begin{qoptions}
\opt{A}{Have two nodal planes}
\opt{B}{Used by sulphur in SF$_6$, for hybridisation}
\opt{C}{Probability of finding electron is maximum along the $x$ and $y$ axis}
\opt{D}{Can accomodate maximum two electrons}
\end{qoptions}
\qn{10} In which of the following arrangements. The order is correct according to the property indicating against it :-
\begin{qoptions}
\opt{A}{Increasing size : Al$^{3+}<$ Mg$^{+2}<$ Na$^+<$ F$^\Theta$}
\opt{B}{Increasing IE$_2$ : C $<$ B $<$ N $<$ O}
\opt{C}{Increasing EN : I $<$ Br $<$ F $<$ Cl}
\opt{D}{Increasing metallic radius : Li $<$ Na $<$ K $<$ Rb}
\end{qoptions}
\sect{SECTION-II(i)}
\paraq{1}{2} Read the following passage and answer the questions based on it.

In a reaction, formation of maximum number of moles of any of the products depends upon the reagent which has stoichiometrically minimum number of moles among the reactants. This reagent, having stoichiometrically minimum number of moles is called Limiting Reagent', which is identified by considering number of moles of all reactants together with stoichiometry (balanced) of the reaction.
\qn{1} 500 mL of 0.4 M BaCl$_2$ solution was mixed with 1000 mL of 0.2 M Na$_3$PO$_4$ solution, then total concentration of ions in resultant solution is \underline{\hspace{1cm}} M? Assume Ba$_3$(PO$_4$)$_2$ completely insoluble in water.
\qn{2} $3.011\times10^{23}$ atoms of Zn, 28 gm Fe and $3.011\times10^{23}$ atoms of sulphur were allowed to react in a container to form Zn(FeS$_2$)$_2$ compound. Find moles of Zn(FeS$_2$)$_2$ is formed ? (Atomic mass of Fe $=56$)
\paraq{3}{4} Hybridisation is the intermixing of atomic orbital of central atom. Hybrid orbital can forms $\sigma$ bond but does not forms $\pi$ bond. Usually but they can have lone pair electron
\qfig{p1_chem_s3_para3_4_1}
\qn{3} If difference of bond angel between BCl$_3$ \& BF$_4^-$ is $x$ and difference between bond angle of BeCl$_2$ \& BeCl$_3^-$ is $y$. Find $(x+y)$.

(Use $\theta=109^\circ28'$ for perfect tetrahedron)
\qn{4} How many statement are incorrect

(1) Silane is formed in the ground state of Si

(2) sp$^2$ hybrid Nitrogen is more E.N. then sp hybrid Nitrogen

(3) Nitrogen atom is sp$^3$ hybridise in N(SiH$_3$)$_3$.

(4) Allene is an example of planer molecule
\paraq{5}{6} According to VSEPR theory electron pairs tend to minimize repulsions and these are in the order -- lone pair--lone pair $>$ lone pair--bond pair $>$ bond pair--bond pair. (Here bond pair refers to a single bond.) shapes of the molecules depend upon repulsions between bond pair and lone pair electrons.

The double bond is in need of more space as compared to the single bond. The repulsion order in relation to the bonds is as follow :

double bond--double bond $>$ double bond--single bond $>$ single bond--single bond

Keeping the central atom (having lone pair) same, if the electronegativity of the surrounding atom increases, the bond angle will decrease provided no other factors like size and back bonding play any role.

Keeping the surrounding atoms same, if the electronegativity of the central atom (having the lone pair) increases, the bond angle increases.
\qn{5} Find out the total number of species having atleast one adjacent bond angle is $\le90^\circ$.

BrF$_5$, XeO$_3$, SiCl$_4$, SOF$_4$, SO$_3$, PCl$_5$, IO$_2$F$_2^-$
\qn{6} Among the following, find out total number of species having planar as well as each bond length is identical.

SF$_4$, ClF$_3$, ICl$_2^-$, I$_3^-$, SCl$_2$, PCl$_5$, IF$_4^-$, BrF$_5$
\sect{SECTION-II(ii)}
\qn{7} 30 ml of a solution containing 9.15 g/L of an oxalate K$_x$H$_y$(C$_2$O$_4$)$_z$.nH$_2$O required for titration 27 ml of 0.12 N NaOH and 36 ml of 0.12 N KMnO$_4$ for oxidation. Find the value of $y$.

($x$, $y$ and $z$ are simplest integer co-efficients)
\qn{8} Number of 3$^\circ$ hydrogens present in the given compound
\qfig{p1_chem_s4_q08_1}
\qn{9} Find the number of chemical species, in which 2p$\pi$--3d$\pi$ back bond is present

B(OH)$_3$, N(SiH$_3$)$_3$, P(CH$_3$)$_3$, N(GeH$_3$)$_3$, BeCl$_2$, O(SiH$_3$)$_2$, OCl$_2$, BCl$_3$, PCl$_3$
\partsec{3}{MATHEMATICS}
\sect{SECTION-I(i)}
\qn{1} The locus of centre of circles which bisect the circumfernce of circle $x^2+y^2=9$ and $x^2+y^2-2x-4y+2=0$ is
\begin{qoptions}
\opt{A}{$2x+4y-1=0$}
\opt{B}{$2x+4y+1=0$}
\opt{C}{$2x-4y+1=0$}
\opt{D}{$2x-4y-1=0$}
\end{qoptions}
\qn{2} If the system of equations $x+4ay+az=0$, $x+3by+bz=0$ and $x+2cy+cz=0$ has a non-trivial solution, then $a,b,c$ are in ($a,b,c$ are distinct non-zero real numbers) -
\begin{qoptions}
\opt{A}{A.P.}
\opt{B}{G.P.}
\opt{C}{H.P.}
\opt{D}{A.G.P.}
\end{qoptions}
\qn{3} Two consecutive sides of a parallelogram are $3x+5y=0$ and $5x+3y=0$. If the equation to one diagonal is $15y+105x=32$ then the equation of the other diagonal is
\begin{qoptions}
\opt{A}{$12y+13x=0$}
\opt{B}{$11y+13x=0$}
\opt{C}{$22y+15x=0$}
\opt{D}{None of these}
\end{qoptions}
\qn{4} Figure shows a circle of unit radius cutting $x$-axis in $A,A'$ and $y$-axis in $B,B'$. $P$ and $Q$ are taken somewhere in second and fourth quadrants respectively and tangents at these points meet in $T$. If lines $PB$, $BT$ and $QB$ meets $x$-axis in $P'$, $T'$ and $Q'$ respectively, then $\dfrac{P'T'}{T'Q'}$ is equal to
\qfig{p1_math_s1_q04_1}
\begin{qoptions}
\opt{A}{4}
\opt{B}{3}
\opt{C}{2}
\opt{D}{1}
\end{qoptions}
\sect{SECTION-I(ii)}
\qn{5} If the pair of lines $x^2+2xy+ay^2=0$ and $ax^2+2xy+y^2=0$ have exactly one line in common and the joint equation of the other two lines is $\alpha x^2+\beta xy+\gamma y^2=0$ where $\alpha,\beta,\gamma\in N$, then lowest value of
\begin{qoptions}
\opt{A}{$\alpha=3$}
\opt{B}{$\beta=9$}
\opt{C}{$\alpha+\gamma=6$}
\opt{D}{$\alpha+\beta=13$}
\end{qoptions}
\qn{6} Let $\alpha,\beta$ and $\gamma$ be the length of sides of a triangle $\Delta PQR$. The possible value of $\left(\dfrac\alpha{\beta+\gamma}+\dfrac\beta{\gamma+\alpha}+\dfrac\gamma{\alpha+\beta}\right)$ can be
\begin{qoptions}
\opt{A}{1.5}
\opt{B}{2}
\opt{C}{3}
\opt{D}{4}
\end{qoptions}
\qn{7} If the system of equations $x-cy-bz=0$, $cx-y+az=0$ and $bx+ay-z=0$ has a non-zero solution and atleast one of $a,b,c$ is a proper fraction, then
\begin{qoptions}
\opt{A}{$a^2+b^2+c^2<3$}
\opt{B}{$abc<-1$}
\opt{C}{$abc>-1$}
\opt{D}{$a^2+b^2+c^2>3$}
\end{qoptions}
\qn{8} The number of solutions of $\sin^5x-\cos^5x=\dfrac1{\cos x}-\dfrac1{\sin x}$ (where $\sin x\ne\cos x$) is :
\begin{qoptions}
\opt{A}{$0$ ; $x\in(0,\pi)-\left\{\dfrac\pi2\right\}$}
\opt{B}{$1$ ; $x\in(0,\pi)-\left\{\dfrac\pi2\right\}$}
\opt{C}{$2$ ; $x\in(0,2\pi)-\left\{\dfrac\pi2,\dfrac{3\pi}2,\pi\right\}$}
\opt{D}{$0$ ; $x\in(0,2\pi)-\left\{\dfrac\pi2,\dfrac{3\pi}2,\pi\right\}$}
\end{qoptions}
\qn{9} If maximum and minimum value of expression $\dfrac{1+\sin\theta}{2-\cos\theta}$ is $M$ and $m$ respectively then :
\begin{qoptions}
\opt{A}{$m=0$}
\opt{B}{$m=\dfrac12$}
\opt{C}{$M=1$}
\opt{D}{$M=\dfrac43$}
\end{qoptions}
\qn{10} A line $L$ is drawn from $P(4,3)$ to meet the lines $L_1\equiv3x+4y+5=0$ and $L_2\equiv3x+4y+15=0$ at points $A$ and $B$ respectively. From point $A$ a line perpendicular to $L$ is drawn meeting the line $L_2$ at $A_1$. Similarly from point $B$, a line perpendicular to $L$ is drawn meeting the line $L_1$ at $B_1$. Thus a parallelogram $AA_1BB_1$ is formed. The equation of $L$ so that area of parallelogram $AA_1BB_1$ is least -
\begin{qoptions}
\opt{A}{$x-7y+17=0$}
\opt{B}{$7x+y-31=0$}
\opt{C}{$x-7y-31=0$}
\opt{D}{$7x+y+17=0$}
\end{qoptions}
\sect{SECTION-II(i)}
\paraq{1}{2} Consider a system of equations :

$x\sin\phi-2y\cos\phi-\lambda z=0$

$x+2y+z=0$

$-x+y+z=0$ (where $\lambda\in R$)
\qn{1} The number of integer values of $\lambda$ for which system have non-trivial solution --
\qn{2} Consider the following statements -

Statement 1 : For $\lambda=1$ there exist $\phi$ for which the system of equation will have infinite solutions.

Statement 2 : For $\lambda=3$ there exist $\phi$ for which the system of equation will have unique solution.

Statement 3 : Number of integral values of $\lambda$ is 3 for the system of equation to have non-trivial solutions.

Then, number of correct statements out of the above, is
\paraq{3}{4} Let $A,B,C$ be 3 sets such that

$A=\left\{(x,y):\dfrac x{\cos\theta}=\dfrac y{\sin\theta}=2,\theta\in R-\left\{\dfrac{n\pi}2\right\},n\in Z\right\}$

$B=\left\{(x,y):\dfrac{x+1}{\cos2205^\circ}=\dfrac y{\sin2205^\circ}=r,r\in R\right\}$

$C=\left\{(x,y):\left(x-\sqrt3\right)^2+(y-1)^2\leqslant a^2:a\in R^+\right\}$

Let $C_1,C_2$ be the curves represented by $A$ and $B$ respectively while $R$ be the region represented by $C$.
\qn{3} A points $P(x_1,y_1)$ and $Q(x_2,y_2)$ be such that $P,Q\in A\cap B$, then point of intersection of tangents drawn to $C_1$ at $P$ and $Q$ is $(a,b)$ then value of $\left(a^2+b^2\right)$ is
\qn{4} If minimum value of `$a$' for which $A\cap C=A$ is $\alpha$, then length of curve $C_2$ lying in $R$ corresponding to $a=\alpha$, is `$k$' then value of $k^2$ is
\paraq{5}{6} As the diagram shows, there is a sequence of the curves $C_0,C_1,C_2$ \ldots\ldots\ldots. It is known that $C_0$ is an regular pentagon having side length unity. We obtain $C_{k+1}$ from $C_k$ by operating as follows : Every side of $C_k$ is divided into five equal parts and then on second and fourth segment of each side an outward regular pentagon is constructed. After constructing such an regular pentagon second and fourth segment from each side is removed.

($k=0,1,2,\ldots\ldots$). Let $P_n$ be the perimeter of $C_n$.
\qfig{p1_math_s3_para5_6_1}
On the basis of above information answer the following :
\qn{5} The number of sides in $C_{11}$ is $q(11^{11})$, where $q$ is equal to
\qn{6} $\displaystyle\sum_{k=0}^{\infty}\dfrac1{P_k}$ equals to
\sect{SECTION-II(ii)}
\qn{7} If
\par\noindent\hfill$\begin{vmatrix}1&1&1\\\cos\alpha&\cos(\alpha+\beta)&\cos(\alpha+\beta+\gamma)\\\sin\alpha&\sin(\alpha+\beta)&\sin(\alpha+\beta+\gamma)\end{vmatrix}=0$\hfill\null\par
such that $\alpha\in[0,2\pi]$ and $\beta,\gamma\in(0,2\pi)$, then $\dfrac{(\beta+\gamma)}{\alpha_{\max}}$ is equal to
\qn{8} From a point `$P$' on the line $2x+y+4=0$ which is nearest circle $x^2+y^2-12y+35=0$, tangents are drawn to the given circle. If the area of quadrilateral formed by these pair of tangents and pair of radii is $k$, then value of $k^2$ is
\qn{9} Let a variable line $\lambda(x-2y+3)+\mu(x+y-6)=0$, $\lambda,\mu\in R$ divides a circle $S=0$ into two regions of equal area. The circle $S=0$ passes through the points on co-ordinate axes where one of the lines cuts off equal non zero intercept on the co-ordinate axes. If area of circle is $k\pi$, then value of $k$ is
\columnbreak
\papersec{2}
\partsec{1}{PHYSICS}
\sect{SECTION-I}
\qn{1} The figure shows an isosceles, uniform, thin triangular plate of mass $M$ and hypotenuse $L$. The angle at the apex is $90^\circ$. The apex lies at the origin and the hypotenuse is parallel to $x$ - axis.
\qfig{p2_phy_s1_q01_1}
\begin{qoptions}
\opt{A}{The moment of inertia of the plate about the $z$ - axis is $\dfrac{ML^2}6$.}
\opt{B}{The moment of inertia of the plate about the $x$ - axis is $\dfrac{ML^2}8$.}
\opt{C}{The moment of inertia of the plate about its hypotenuse axis is $\dfrac{ML^2}{24}$.}
\opt{D}{The moment of inertia of the plate about the $y$ - axis is $\dfrac{ML^2}{24}$.}
\end{qoptions}
\qn{2} A sphere of radius $R$ is supported by a rope attached to the wall as shown in the figure. The point of suspension $p$ is at a distance $\dfrac{3R}2$ from the point of contact of the sphere with the wall. Coefficient of static friction between the sphere and the wall is $\mu$, $\theta$ is angle made by the string with the wall and the sphere is in equilibrium. The correct option(s) of the following is/are.
\qfig{p2_phy_s1_q02_1}
\begin{qoptions}
\opt{A}{If $\theta=\tan^{-1}\left(\dfrac23\right)$, force of friction on the sphere is $\dfrac{2\mu mg}3$.}
\opt{B}{If $\theta=\tan^{-1}\left(\dfrac23\right)$, force of friction on the sphere is 0.}
\opt{C}{If $\theta=45^\circ$, the minimum coefficient of friction between sphere and the wall is 0.5 for no slipping at the point of contact.}
\opt{D}{If $\theta=45^\circ$, force of friction on the sphere is zero.}
\end{qoptions}
\qn{3} The net external torque on a system of particles about an axis is zero. Which of the following are compatible with it?
\begin{qoptions}
\opt{A}{The forces may be acting radially from a point on the axis.}
\opt{B}{The forces may be acting on the axis of rotation.}
\opt{C}{The forces may be acting parallel to the axis of rotation.}
\opt{D}{The torque caused by some forces may be equal and opposite to that caused by other forces.}
\end{qoptions}
\qn{4} A particle is suspended from a rigid support by a massless inextensible string. The particle rotates in a horizontal circle with constant angular speed as in a conical pendulum. The correct option(s) of the following is/are
\begin{qoptions}
\opt{A}{Net torque on the particle about the axis of rotation is zero.}
\opt{B}{Torque due to tension in the string about the axis of rotation is zero.}
\opt{C}{Torque due to weight of the particle about the point of suspension changes with time.}
\opt{D}{Torque due to weight of the particle about the point of the suspension does not change with time.}
\end{qoptions}
\qn{5} Particles $A$ and $B$ are connected by a massless, inextensible cord which goes around the fixed massless pulley. If the system is released from rest with the spring initially unstretched, then:
\qfig{p2_phy_s1_q05_1}
\begin{qoptions}
\opt{A}{The maximum extension in the spring is $\dfrac{2mg}K$.}
\opt{B}{The maximum velocity of particle $C$ is $\sqrt{\dfrac{3m}{4K}}g$.}
\opt{C}{The maximum extension in the spring is $\dfrac{mg}K$.}
\opt{D}{The maximum velocity of particle $C$ is $\sqrt{\dfrac{3m}{2K}}g$.}
\end{qoptions}
\qn{6} In the given figure, a long platform of mass $m$ is placed on frictionless surface. Two blocks of masses $4m$ and $m$ ($m=10$kg) are placed on the platform. For both blocks, the coefficient of static friction with platform equal to 0.16 and the coefficient of kinetic friction equal to 0.10. The blocks are connected by a light ideal string through a light pulley (mounted at a movable massless stand ) which is acted upon by an unknown horizontal force $F$. The acceleration of the platform is 2 m/s$^2$. Which of the following options is/are correct? (Take : $g=10$ m/s$^2$)
\qfig{p2_phy_s1_q06_1}
\begin{qoptions}
\opt{A}{$F$ is 180 N.}
\opt{B}{Acceleration of the $4m$ block is 8 m/s$^2$.}
\opt{C}{Acceleration of the $m$ block is 2 m/s$^2$.}
\opt{D}{The sum of magnitudes of accelerations of the two blocks is 10 m/s$^2$.}
\end{qoptions}
\sect{SECTION-II}
\qn{1} In the diagram shown, there is no friction at any contact surface. Initially, the spring has no deformation. What will be the maximum deformation in the spring if a constant horizontal force $F$ is applied as shown? Consider all the strings to be sufficiently large. The spring constant is $K$. If the maximum deformation in the spring is $QF/K$, find $Q$.
\qfig{p2_phy_s2_q01_1}
\qn{2} Three small balls of equal mass($m$) are suspended from a thread and two springs of same force constant ($k$) such that the distance between the first and the second ball and the second and third ball are equal. Thus the centre of mass of the whole system coincides with the second ball. The thread supporting the upper ball is cut and system starts a free fall. Then the distance of the centre of mass of the system from the second ball when both the springs acquire their natural length in the falling system is $Nmg/k$. Find $N$.
\qfig{p2_phy_s2_q02_1}
\qn{3} Both the blocks shown in the given arrangement are given together a horizontal velocity towards right. If $x$ ms$^{-2}$ be the subsequent acceleration of centre of mass of the system of blocks, find the value of $x$ (in ms$^{-2}$). Given, $m_A=2m_B$ and $g=10$ ms$^{-2}$, all in SI units.
\qfig{p2_phy_s2_q03_1}
\qn{4} A small particle of mass $m$ starts sliding down from rest along the smooth surface of a fixed hollow hemisphere of mass ($M=4m$). The distance of centre of mass of (particle $+$ hemisphere) from centre $O$ of hemisphere, when the particle separates from the surface of hemisphere is ($R/\alpha$). Find the value of $\alpha$.
\qfig{p2_phy_s2_q04_1}
\qn{5} For the figure shown, the wedge $+$ particle is resting at a vertical wall. The particle of mass $\left(m=\dfrac M2\right)$ is released from rest from the position shown. The wedge has mass $M$. The maximum height reached by the particle on the circular groove is ($R/\alpha$). Find the value of $\alpha$.
\qfig{p2_phy_s2_q05_1}
\qn{6} In the $x$-$y$ plane there exist a closed figure of total area 25 cm$^2$. Find the magnitude of work done by the $\vec F=K\left(y\hat i-x\hat j\right)$ in one complete round around the circumference of the figure. Take $K=100$N/m.
\qfig{p2_phy_s2_q06_1}
\sect{SECTION-III}
\qn{1} In the figure shown $ABCD$ is a uniform, thin plate of mass 12 kg. The moment of inertia of this lamina about an axis perpendicular to plane of lamina and passing through point $D$ is $64\eta$ kgm$^2$ then $\eta=$ ?
\qfig{p2_phy_s3_q01_1}
\qn{2} The uniform rod $AB$ of mass $m$ and length $\ell$ is released from rest from vertical dotted position as shown in the figure. The normal force as a function of angle $\theta$ turned by the rod, which is exerted on the rod by the rough horizontal ground as it rotates in a vertical plane about its end $A$ in contact with ground is $mg\left(\dfrac{3\cos\theta-1}\eta\right)^2$ then $\eta=$ ? (assume friction due to ground is sufficient to prevent skidding of end $A$ of rod in contact with the ground)
\qfig{p2_phy_s3_q02_1}
\qn{3} Two identical particles $A$ and $B$ connected by an inextensible, ideal, massless, elastic rope of length $2\ell$ are placed at rest on a smooth horizontal floor at a separation of $\sqrt3\ell$. At $t=0$, $A$ is given a horizontal velocity $u$ perpendicular to the line joining $AB$ as shown in the figure. If the distance moved by $A$ till the time $t=\dfrac{11\ell}u$ is $\sqrt3\left(2+\sqrt x\right)\ell$, where $x$ is a positive integer, find the value of $x$.
\qfig{p2_phy_s3_q03_1}
\qn{4} A hemisphere of mass $2m$ is free to slide with its base on a smooth horizontal table. A particle of mass $m$ is placed on the hemisphere at an angular distance $\alpha=\dfrac\pi3$ from the top of the hemisphere. The radius to the point of contact at which the particle leaves the surface makes an angle $\theta$ with the vertical given by the equation $\cos3\theta+3a\cos\theta+b=0$, where $a$ and $b$ are integers. Find $a+b$.
\qn{5} A boy and a man carry a uniform bar of length $L$ horizontally so that the boy gets $\dfrac13$ rd of the weight of the bar. If the boy is at one end of the bar, the distance of the man from the other end is $\dfrac Ln$. The value of $n$ is
\qn{6} A bead mass 2 kg is fitted onto a rod with a length of 8 m and can move on it with friction having the coefficient of friction 0.5. At the initial moment the bead is in the middle of the rod. The rod moves translational in a horizontal plane with an acceleration 1 m/s$^2$ in the direction forming an angle $37^\circ$ with the rod in gravity free region. Find the time (in sec) when the bead will leave the rod.
\qfig{p2_phy_s3_q06_1}
\partsec{2}{CHEMISTRY}
\sect{SECTION-I}
\qn{1} Graph between $\log(1/n)$ v/s $\log P$ is plotted for an ideal gas enclosed (as shown in diagram) 24.63 litre container at three different temperatures.
\qfig{p2_chem_s1_q01_1}
If $T_1=\dfrac{T_2}4=3T_3$ (where $T_1,T_2,T_3$ are temperature in kelvin of graph 1,2 and 3) then

($R=0.0821$ L-atm/K-mol $\log2=0.301$)
\begin{qoptions}
\opt{A}{$T_1$, $T_2$ and $T_3$ are 300, 1200 and 100 K respectively.}
\opt{B}{Slope of graphs 1, 2 and 3 are different.}
\opt{C}{Intercept of graph 1, 2 and 3 is same.}
\opt{D}{Intercept of graph 2 is 0.602.}
\end{qoptions}
\qn{2} A mixture of C$_2$H$_2$ and C$_3$H$_8$ including sufficient O$_2$ occupied a certain volume at 80 mm Hg. The mixture was completely burnt to CO$_2$ and H$_2$O($\ell$) and observed pressure of CO$_2$ was found to be 42.5 mm Hg at the same temperature and volume. Select the correct statement(s), neglecting the vapour pressure of water :
\begin{qoptions}
\opt{A}{The final pressure of reaction mixture was 42.5 mmHg.}
\opt{B}{Pressure of oxygen in the initial mixture is 62.5 mmHg}
\opt{C}{Pressure of C$_2$H$_2$ in the initial mixture is 32 mmHg.}
\opt{D}{Moles fraction of C$_3$H$_8$ in the initial mixture is 0.6.}
\end{qoptions}
\qn{3} In which of the following ions, positive charge / vacant orbital is in resonance?
\begin{qoptions}
\opt{A}{\qfigi{p2_chem_s1_q03_A}}
\opt{B}{\qfigi{p2_chem_s1_q03_B}}
\opt{C}{\qfigi{p2_chem_s1_q03_C}}
\opt{D}{\qfigi{p2_chem_s1_q03_D}}
\end{qoptions}
\qn{4} Statements which are correct about cubane
\begin{qoptions}
\opt{A}{D.U. $=6$}
\opt{B}{D.U. $=5$}
\opt{C}{No. of 3$^\circ$ H $=8$}
\opt{D}{No. of 2$^\circ$C $=4$}
\end{qoptions}
\qn{5} Select the correct statement(s) about the compound NO[BF$_4$] :

[All atoms have complete octet in given species]
\begin{qoptions}
\opt{A}{If has 5$\sigma$ and 2$\pi$ bond}
\opt{B}{Formal charge on B is $+1$ and on oxygen is $-1$}
\opt{C}{Molecule has even number of valence electrons}
\opt{D}{Hybridisation of B in NO[BF$_4$] is sp$^3$}
\end{qoptions}
\qn{6} Which of the following order is/are \textbf{INCORRECT} for bond angle
\begin{qoptions}
\opt{A}{NH$_3<$ PH$_3<$ AsH$_3<$ SbH$_3$}
\opt{B}{OF$_2<$ OCl$_2<$ OBr$_2>$ H$_2$O}
\opt{C}{BF$_3<$ BCl$_3<$ BBr$_3<$ BI$_3$}
\opt{D}{NO$_2^+>$ NO$_2^-<$ NO$_3^-$}
\end{qoptions}
\sect{SECTION-II}
\qn{1} What would be the maximum volume (in litre) of 3M HCl solution that can be prepared by using 2M, 1L HCl and 5M, 2L HCl kept in seperate vessel.
\qn{2} In a reaction, P$_4$ is converted into H$_3$PO$_4$. one mole of P$_4$ will contain how many equivalents of it in this reaction.
\qn{3} How many of these compounds / ions are aromatic?
\qfig{p2_chem_s2_q03_1}
\qfig{p2_chem_s2_q03_2}
\qfig{p2_chem_s2_q03_3}
\qfullbegin
\qn{4} Find out the number of compounds / ions in which hyper conjugation is present?
\qfig{p2_chem_s2_q04_1}
\qfullsplit
\qfig{p2_chem_s2_q04_2}
\qfig{p2_chem_s2_q04_3}
\qfullend
\qn{5} If $x$ is the total number of p$\pi$-p$\pi$ bond present and $y$ is the total number of d$\pi$-p$\pi$ bond present then find the value of $\dfrac{2(x+y)}x$ from the following molecules.

P$_4$O$_{10}$, XeO$_3$, SO$_3$, SO$_2$Cl$_2$, HClO$_4$, SO$_2$
\qn{6} If electronegativity difference between two atoms $A$ and $B$ is 2.0, then the percentage of covalent character in the molecule is :
\sect{SECTION-III}
\qn{1} A solution of fatty acid $X$ (GMM $=252$) in benzene contain 4.2gm of acid per litre. When this solution is dropped in a water surface the benzene evaporates and fatty acid $X$ forms a mono layer film of solid type. If we wish to cover 10 m$^2$ area with monolayer film of fatty acid $X$, what volume of solution (in ml) should we use. Each fatty acid molecule occupy 0.2 nm$^2$ area. $N_A=6\times10^{23}$.
\qn{2} The formula weight of an acid is 82. In a titration, 100 cm$^3$ of a solution of this acid containing 39 g of the acid per litre were completely neutralised by 95.0 cm$^3$ of aqueous NaOH containg 40.0g of NaOH per litre. Then the basicity of the acid is \underline{\hspace{1cm}}.

(Answer to nearest integer)
\qn{3} Number of F.G. present in compound H$_2$N--CHO are $x$ and number of F.G. present in compound HO--CHO are $y$ then calculate the value of $(x+y)$?
\qn{4} How many resonating structure are possible in
\qfig{p2_chem_s3_q04_1}
\qn{5} Find the total number of correct orders among the following.

(1) Order of electron affinity : S $>$ O$^-$

(2) Order of ionic radius : Al$^{3+}<$ Mg$^{2+}<$ O$^{2-}$

(3) Order of Electronegativity : Si $<$ P $<$ C $<$ N

(4) Order of atomic radius : O $<$ S $<$ Se

(5) Order of second ionization energy : C $<$ N $<$ F $<$ O

(6) Order of ionic radii : S$^{2-}<$ Cl$^-<$ K$^+<$ Ca$^{2+}$
\qn{6} How many of the following can react with alkali :

Mn$_2$O$_7$, ZnO, Al$_2$O$_3$, BeO, SeO$_2$, Cl$_2$O$_7$, Fe$_2$O$_3$, As$_2$O$_3$
\columnbreak
\partsec{3}{MATHEMATICS}
\sect{SECTION-I}
\qn{1} If $T_n=\begin{vmatrix}n-1&n-2&n-6\\2n-4&2n-6&2n-11\\3n-9&3n-12&3n-18\end{vmatrix}$ where $n\in N$, then which of the following is/are true ?
\begin{qoptions}
\opt{A}{$\displaystyle\prod_{n=1}^{m}T_n=6^m$}
\opt{B}{$\displaystyle\sum_{n=1}^{10}T_n=60$}
\opt{C}{$\dfrac{T_{n+1}}{T_n}=2$}
\opt{D}{$T_{100}.T_{101}=48$}
\end{qoptions}
\qn{2} Let $A,B$ and $C$ are three points on a circle such that arc $AB=6$, arc $BC=8$, arc $CA=10$, then
\begin{qoptions}
\opt{A}{$\Delta ABC$ is acute angled}
\opt{B}{$\angle B,\angle A,\angle C$ are in A.P}
\opt{C}{$\angle A,\angle B,\angle C$ are in A.P}
\opt{D}{Area of $\Delta ABC$ is $\dfrac{36\left(3+\sqrt3\right)}{\pi^2}$}
\end{qoptions}
\qn{3} The equation of the straight line equally inclined to the coordinate axes and equidistant from the points $(2,0)$ and $(0,4)$ can be-
\begin{qoptions}
\opt{A}{$y=x-1$}
\opt{B}{$y=x+1$}
\opt{C}{$x+y=3$}
\opt{D}{$x+y+3=0$}
\end{qoptions}
\qn{4} In $\Delta ABC$, equation of side $BC$ is $x+y=6$, the circumcentre and orthocentre are $(3,1)$ and $(2,2)$ respectively. Then
\begin{qoptions}
\opt{A}{reflection of orthocentre about side $BC$ is $(4,4)$}
\opt{B}{circumradius of $\Delta ABC$ is $\sqrt{10}$}
\opt{C}{circumradius of $\Delta ABC$ is $\dfrac{\sqrt{10}}2$}
\opt{D}{centroid of $\Delta ABC$ is $\left(\dfrac83,\dfrac43\right)$}
\end{qoptions}
\qn{5} Consider the curve
\par\noindent\hfill$\begin{aligned}S&=(\sin\theta)x^2+2(\cos\theta)xy+(\sin\theta)y^2\\&\quad+2(\cos\theta)x+2(\cos\theta)y+\sin\theta=0\end{aligned}$\hfill\null\par
Now, choose the correct statement(s)
\begin{qoptions}
\opt{A}{If $\theta=\dfrac{3\pi}4$, then $S=0$ represents a pair of straight lines.}
\opt{B}{If $S=0$ represents a pair of real straight lines, then $\theta=n\pi+\dfrac\pi4$, $n\in Z$}
\opt{C}{If $S=0$ represents a pair of real straight lines, then they are parallel.}
\opt{D}{There exist some $\theta$ for which $S=0$ represents no real locus.}
\end{qoptions}
\qn{6} If $4a^2-5b^2+6a+1=0$ and the line $ax+by+1=0$ touches a fixed circle, then (where $a,b\in R$)
\begin{qoptions}
\opt{A}{centre of circle is $(3,0)$}
\opt{B}{radius of circle $=\sqrt5$}
\opt{C}{circle passes through $(0,1)$}
\opt{D}{None of these}
\end{qoptions}
\sect{SECTION-II}
\qn{1} If $\left[\displaystyle\sum_{r=0}^{n}2^r\tan\left(2^r\theta\right)\right]+2^{n+1}\cot\left(2^{n+1}\theta\right)=k\left(\cot\dfrac\theta2-\tan\dfrac\theta2\right)$ is valid for all $\theta$ (wherever defined), then $k$ is equal to
\qn{2} Consider a sequence $\{a_n\}$ with $a_1=2$ and $a_n=\dfrac{a_{n-1}^2}{a_{n-2}}$ (for all $n\ge3$, terms of the sequence being distinct). Given that $a_2$ and $a_5$ are positive integers and $a_5\le162$, then the smallest possible value(s) of $a_5$ can be
\qn{3} If points $(0,0)$, $(\lambda,0)$ $(0,\lambda)$ and $(\mu-1,\mu+3)$ are always concyclic $\forall\ \mu\in R$ (where $\lambda\in R$ and $\lambda\ne0$), then minimum value of $|\lambda|$ is
\qn{4} A circle $x^2+y^2+4x-2\sqrt2\,y+c=0$ is the director circle of circle $S_1$ and $S_1$ is the director circle of circle $S_2$ and so on. If the sum of radii of all these circles is 2, then the value of $c$ is equal to $\sqrt n$ where $n\in N$.

Find the value of $n$.
\qn{5} Let $L_1$, $L_2$ and $L_3$ be the length of tangents drawn from a point $P$ to the circles $x^2+y^2=4$, $x^2+y^2-4x=0$ and $x^2+y^2-4y=0$ respectively. If $L_1^4=L_2^2L_3^2+16$, then locus of $P$ are the curves $C_1$(a straight line) and $C_2$(a circle). Let the circumcentre of the triangle formed by $C_1$ and other two lines which are inclined at an angle of $45^\circ$ with $C_1$ and is tangent to $C_2$ is $(\alpha,\beta)$, then $(\alpha+2\beta)$ equals
\qn{6} The co-ordinates of a point $P(p,q)$ on the line $4x-y=10$ such that $|PA-PB|$ is maximum where $A$ is $(2,3)$ and $B$ is $(1,-2)$, then the absolute value of $\left(\dfrac qp\right)$ is
\sect{SECTION-III}
\qn{1} If point $(2a,a)$ lies inside the parallelogram formed by the lines $2x-y+1=0$, $2x-y-20=0$, $x+y-1=0$ and $x+y-8=0$, then sum of all integral values of `$a$' is -
\qn{2} Consider a curve $|x|+|y|=1$ such that lines $y=mx$, $y=nx$ make points of intersection in the same quadrant. Let $A$, $B$ be two such points lying in the same quadrant sucht that $OA$, $AB$, $OB$ are in A.P. If $\theta$ is angle between $OA$ and $OB$ then the maximum angle $\theta$ is $\dfrac\pi k$ ($O$ is origin) , where $k$ is equal to
\qn{3} If $\alpha,\beta,\gamma$ are the roots of $x^3+x^2-2x+1=0$, then the value of $\sqrt\Delta$ is equal to \underline{\hspace{1.5cm}}

where $\Delta=\begin{vmatrix}\gamma^2&\beta^2&2\beta\gamma-\alpha^2\\2\alpha\gamma-\beta^2&\alpha^2&\gamma^2\\\alpha^2&2\alpha\beta-\gamma^2&\beta^2\end{vmatrix}$
\qn{4} Consider a triangle $ABC$ with $A(0,0)$, $B\left(2\sqrt3,0\right)$ and $C\left(\sqrt3,3\right)$. If $L$ is locus of the points from which tangents of equal lengths can be drawn to its incircle \& excircle opposite to vertex $A$, then the minimum distance of $L$ from origin.
\qn{5} Let a circle with centre $O$ and radius $k$. Three points $P$, $Q$, $R$ are chosen on a circle such that if length of chord $PQ$, $QR$, $RP$ are $x$, $y$ \& $z$ respectively, $x^2+y^2+z^2=5$, $\angle QOR=60$ and Area of $\Delta PQR$ is $\dfrac{\sqrt3}4$ unit, then the value of radius $k$ of the circle is equal to
\qn{6} Consider a family of lines $(2k+1)x+(3k-2)y+(3-8k)=0$ where $k\in R$. The minimum area of the triangle formed by a line of given family with positive semi axes is equal to
\end{multicols}
\end{document}
